\heading{实验物理中的统计方法-特征函数}





\begin{question}
证明高斯分布

\[
f(x;\mu,\sigma^2)=\frac{1}{\sqrt{2\pi\sigma^2}}\exp\left[-\frac{(x-\mu)^2}{2\sigma^2}\right]
\]

的特征函数为

\[
\phi(k)=\exp\left(i\mu k-\frac12\sigma^2k^2\right).
\]
\begin{solution}
设

\[
f(x;\mu,\sigma^2)=\frac{1}{\sqrt{2\pi}\sigma}\exp\!\left[-\frac{(x-\mu)^2}{2\sigma^2}\right].
\]

则

\[
\phi(k)=\int_{-\infty}^{\infty} e^{ikx} f(x;\mu,\sigma^2)\,dx.
\]

把指数配方：

\[
ikx-\frac{(x-\mu)^2}{2\sigma^2}
=
-\frac{(x-\mu-i\sigma^2 k)^2}{2\sigma^2}+ik\mu-\frac12\sigma^2k^2.
\]

因此高斯积分部分仍为 1，得到

\[
\phi(k)=\exp\!\left(ik\mu-\frac12\sigma^2k^2\right).
\]
\end{solution}
\end{question}



\begin{question}
证明指数分布

\[
f(x;\xi)=\frac1\xi e^{-x/\xi}
\]

的特征函数为

\[
\phi(k)=\frac{1}{1-ik\xi}.
\]
\begin{solution}
设

\[
f(x;\xi)=\frac1{\xi}e^{-x/\xi},\qquad x\ge 0.
\]

则

\[
\phi(k)=\int_0^\infty e^{ikx}\frac1{\xi}e^{-x/\xi}\,dx
=\frac1{\xi}\int_0^\infty e^{-x(1/\xi-ik)}\,dx.
\]

只要 \(\Re(1/\xi-ik)=1/\xi>0\)，积分收敛，故

\[
\phi(k)=\frac1{\xi}\cdot \frac1{1/\xi-ik}
=\frac1{1-ik\xi}.
\]
\end{solution}
\end{question}



\begin{question}
证明自由度为 \(n\) 的 \(\chi^2\) 分布

\[
f(z;n)=\frac{1}{2^{n/2}\Gamma(n/2)}z^{n/2-1}e^{-z/2}
\]

的特征函数为

\[
\phi(k)=(1-2ik)^{-n/2}.
\]

提示：证明中需要用到伽马函数，定义为

\[
\Gamma(x)=\int_0^\infty e^{-t}t^{x-1}\,dt.
\]
\begin{solution}
若 \(Z\sim \chi^2_n\)，其密度为

\[
f(z;n)=\frac{1}{2^{n/2}\Gamma(n/2)}z^{n/2-1}e^{-z/2},\qquad z>0.
\]

则

\[
\phi(k)=\frac{1}{2^{n/2}\Gamma(n/2)}
\int_0^\infty z^{n/2-1}e^{-z(1/2-ik)}\,dz.
\]

利用伽马积分 \[
\int_0^\infty z^{a-1}e^{-bz}\,dz=\frac{\Gamma(a)}{b^a},
\] 得

\[
\phi(k)=\frac{1}{2^{n/2}}\frac{1}{(1/2-ik)^{n/2}}
=(1-2ik)^{-n/2}.
\]

故

\[
\phi_{\chi^2_n}(k)=(1-2ik)^{-n/2}.
\]
\end{solution}
\end{question}



\begin{question}
假定随机变量 \(x_1,\ldots,x_n\) 相互独立，且都服从均值为 \(\mu\)、方差为
\(\sigma^2\) 的高斯分布。第 5 章和第 6 章给出，可以用样本均值

\[
\bar x=\frac1n\sum_{i=1}^{n}x_i
\]

作为均值 \(\mu\) 的估计量。

\begin{enumerate}
\item 求样本均值的特征函数。

\item 利用特征函数证明 \(\bar x\) 服从高斯分布，并求其均值和方差。
\end{enumerate}
\begin{solution}
设 \(x_1,\dots,x_n\) 独立同分布，且

\[
x_i\sim N(\mu,\sigma^2),\qquad \bar x=\frac1n\sum_{i=1}^n x_i.
\]

由独立性，

\[
\phi_{\bar x}(k)
=E\!\left[\exp\!\left(ik\frac1n\sum_i x_i\right)\right]
=\prod_{i=1}^n \phi_{x_i}(k/n).
\]

代入 10.1 的结果：

\[
\phi_{\bar x}(k)
=\left[\exp\!\left(i\mu\frac{k}{n}-\frac12\sigma^2\frac{k^2}{n^2}\right)\right]^n
=\exp\!\left(ik\mu-\frac12\frac{\sigma^2}{n}k^2\right).
\]

这正是高斯分布 \(N(\mu,\sigma^2/n)\) 的特征函数，所以

\[
\bar x\sim N\!\left(\mu,\frac{\sigma^2}{n}\right).
\]
\end{solution}
\end{question}



\begin{question}
利用特征函数证明高斯分布的前 4 阶代数矩为

\[
\begin{split}
E[x]&=\mu,\\
E[x^2]&=\mu^2+\sigma^2,\\
E[x^3]&=\mu^3+3\mu\sigma^2,\\
E[x^4]&=3(\mu^2+\sigma^2)^2-2\mu^4.
\end{split}
\]
\begin{solution}
由

\[
\phi(k)=\exp\!\left(ik\mu-\frac12\sigma^2k^2\right)
\]

在 \(k=0\) 展开，或者直接使用 \[
E[X^m]=\left.\frac{1}{i^m}\frac{d^m\phi}{dk^m}\right|_{k=0},
\] 可得

\[
E[X]=\mu,
\] \[
E[X^2]=\mu^2+\sigma^2,
\] \[
E[X^3]=\mu^3+3\mu\sigma^2,
\] \[
E[X^4]=\mu^4+6\mu^2\sigma^2+3\sigma^4
=3(\mu^2+\sigma^2)^2-2\mu^4.
\]

故书上给出的结果成立。
\end{solution}
\end{question}



\begin{question}
\begin{enumerate}
\item 利用特征函数证明自由度为 \(n\) 的 \(\chi^2\) 分布的均值和方差分别为
\(n\) 和 \(2n\)。

\item 假设 \(z\) 服从自由度为 \(n\) 的 \(\chi^2\) 分布。证明：在大 \(n\)
极限下，\(\chi^2\) 分布变为均值为 \(\mu=n\)、方差为 \(\sigma^2=2n\)
的高斯分布。提示：证明中需要定义变量

\[
y=\frac{z-n}{\sqrt{2n}},
\]

\(y\) 的均值为零，标准差为 1。证明 \(y\) 的特征函数为

\[
\phi_y(k)=e^{-ik\sqrt{n/2}}\phi_z\left(\frac{k}{\sqrt{2n}}\right).
\]

将 \(\phi_y(k)\) 的对数展开，并只保留大 \(n\)
极限下不消失的项，然后再变换回变量 \(z\) 得到要证明的结果。
\end{enumerate}
\begin{solution}
\begin{enumerate}
\item 均值与方差

由

\[
\phi(k)=(1-2ik)^{-n/2}
\]

取对数： \[
\log \phi(k)=-\frac n2 \log(1-2ik).
\]

也可直接求矩母函数 \(M(t)=(1-2t)^{-n/2}\)。于是

\[
E[Z]=M'(0)=n,\qquad \mathrm{Var}(Z)=M''(0)-M'(0)^2=2n.
\]

故

\[
E[Z]=n,\qquad \mathrm{Var}(Z)=2n.
\]

\item 渐近高斯性

定义标准化变量

\[
Y=\frac{Z-n}{\sqrt{2n}}.
\]

则其特征函数为

\[
\phi_Y(k)=e^{-ik\sqrt{n/2}}\phi_Z\!\left(\frac{k}{\sqrt{2n}}\right)
=e^{-ik\sqrt{n/2}}\left(1-\frac{2ik}{\sqrt{2n}}\right)^{-n/2}.
\]

取对数并在 \(n\to\infty\) 下展开：

\[
\log\phi_Y(k)
=
-ik\sqrt{\frac n2}
-\frac n2 \log\!\left(1-\frac{\sqrt{2}ik}{\sqrt n}\right)
\to -\frac12 k^2.
\]

故 \[
\phi_Y(k)\to e^{-k^2/2},
\] 这就是标准正态分布的特征函数，因此

\[
\frac{Z-n}{\sqrt{2n}}\xrightarrow[n\to\infty]{d}N(0,1).
\]
\end{enumerate}
\end{solution}
\end{question}



\begin{question}
假设 \(n\) 个相互独立的随机变量 \(x_1,\ldots,x_n\)
都服从标准高斯分布，即对所有 \(i=1,\ldots,n\)，

\[
\phi(x_i)=\frac{1}{\sqrt{2\pi}}e^{-x_i^2/2}.
\]

考虑下面这个变量的性质：

\[
y=\left(\sum_{i=1}^{n}x_i^2\right)^{1/2}.
\]

\begin{enumerate}
\item 首先只考虑某一个 \(x_i\)。通过变量变换，证明 \(u=x_i^2\)
的概率密度函数为

\[
f(u)=\frac{1}{\sqrt{2\pi u}}e^{-u/2}.
\]

这是自由度为 1 的 \(\chi^2\) 分布。

\item 证明 \(u\) 的特征函数为

\[
\phi_u(k)=\frac{1}{\sqrt{1-2ik}}.
\]

\item 利用相加定理，求下面变量的特征函数

\[
v=\sum_{i=1}^{n}x_i^2.
\]

\item 利用变量变换，证明 \(y=\left(\sum_{i=1}^{n}x_i^2\right)^{1/2}\)
的概率密度函数为

\[
h(y)=\frac{1}{2^{n/2-1}\Gamma(n/2)}y^{n-1}e^{-y^2/2}.
\]

\item 写出 \(n=3\) 时的概率密度函数。这是 Maxwell-Boltzmann
分布。假设气体中分子的速度分量 \(v_x\)，\(v_y\) 和 \(v_z\)
都服从均值为零、标准差为 \(\sigma\) 的高斯分布。写出分子速度
\(v=(v_x^2+v_y^2+v_z^2)^{1/2}\) 的概率密度函数。

\item 写出 \(n=1\) 时的概率密度函数。即，如果 \(x\) 服从标准高斯分布，则
\(y=|x|\) 的概率密度函数是什么？
\end{enumerate}
\begin{solution}
设 \(x_i\sim N(0,1)\) 相互独立。

\begin{enumerate}
\item 令 \(u=x_i^2\)。由两个反函数分支 \(x_i=\pm\sqrt u\) 得
\[
f_U(u)=\frac{1}{\sqrt{2\pi u}}e^{-u/2},\qquad u>0,
\]
这就是自由度为 1 的 \(\chi^2\) 分布。

\item 因此
\[
\phi_u(k)=E(e^{iku})=(1-2ik)^{-1/2}.
\]

\item 令 \(v=\sum_{i=1}^n x_i^2\)。由独立性，
\[
\phi_v(k)=\prod_{i=1}^n\phi_u(k)=(1-2ik)^{-n/2},
\]
所以 \(v\sim\chi_n^2\)。

\item 因 \(y=\sqrt v\)，即 \(v=y^2\)，有 \(dv/dy=2y\)。于是
\[
h(y)=f_v(y^2)\,2y
=\frac{1}{2^{n/2-1}\Gamma(n/2)}y^{n-1}e^{-y^2/2},
\qquad y>0.
\]

\item 当 \(n=3\) 时，
\[
h(y)=\sqrt{\frac{2}{\pi}}\,y^2e^{-y^2/2},\qquad y>0.
\]
若速度分量 \(v_x,v_y,v_z\sim N(0,\sigma^2)\)，则速率 \(r\) 的密度为
\[
f(r)=\sqrt{\frac{2}{\pi}}\frac{r^2}{\sigma^3}e^{-r^2/(2\sigma^2)},\qquad r>0.
\]

\item 当 \(n=1\) 时，
\[
h(y)=\sqrt{\frac{2}{\pi}}e^{-y^2/2},\qquad y>0.
\]
一般地，若 \(x\sim N(0,\sigma^2)\) 且 \(y=|x|\)，则
\[
h(y)=\sqrt{\frac{2}{\pi}}\frac1{\sigma}e^{-y^2/(2\sigma^2)},\qquad y>0.
\]
\end{enumerate}
\end{solution}
\end{question}



\begin{question}
考虑服从柯西分布的随机变量 \(x\)，

\[
f(x)=\frac{1}{\pi}\frac{1}{1+x^2}.
\]

\begin{enumerate}
\item 证明其特征函数为

\[
\phi(k)=e^{-|k|}.
\]

（利用留数定理，\(k>0\) 时选取上半平面的路径，\(k<0\)
时选取下半平面的路径。）

\item 考虑柯西随机变量 \(x\) 的某个样本，样本容量为
\(n\)。利用（a）中得到的特征函数并应用相加定理，证明样本均值
\(\bar x=\frac1n\sum_{i=1}^{n}x_i\) 也服从柯西分布。这是一个特例，即
\(\bar x\)
的概率密度函数不随样本容量的增加而改变，这与柯西分布的各阶矩不存在有关。
\end{enumerate}
\begin{solution}
设

\[
f(x)=\frac{1}{\pi(1+x^2)}.
\]

\begin{enumerate}
\item 特征函数

计算

\[
\phi(k)=\int_{-\infty}^{\infty}\frac{e^{ikx}}{\pi(1+x^2)}\,dx.
\]

利用留数定理：当 \(k>0\) 时取上半平面闭路，包围极点 \(x=i\)；当 \(k<0\)
时取下半平面闭路，包围极点 \(x=-i\)。两种情形都得到

\[
\phi(k)=e^{-|k|}.
\]

\item 样本均值仍服从柯西分布

若 \(x_1,\dots,x_n\) 独立同分布为柯西，则

\[
\phi_{\bar x}(k)
=\prod_{i=1}^n \phi_x(k/n)
=\left(e^{-|k|/n}\right)^n
=e^{-|k|}.
\]

因此 \(\bar x\)
与原变量同分布，仍是柯西分布。这解释了为什么柯西分布没有通常意义上的大数定律行为。
\end{enumerate}
\end{solution}
\end{question}



\begin{question}
狄拉克 \(\delta\) 函数

\[
f(x;\mu)=\delta(x-\mu)
\]

定义为

\[
\delta(x-\mu)=0,\quad x\ne\mu,
\]

\[
\int_{-\infty}^{\infty}\delta(x-\mu)\,dx=1.
\]

即，\(\delta(x-\mu)\) 在 \(x=\mu\)
处为无限尖锐的峰，但在其它地方都等于零。求 \(\delta(x-\mu)\)
的特征函数，并据此得到 \(\delta\) 函数的积分表示。
\begin{solution}
\[
f(x;\mu)=\delta(x-\mu).
\]

其特征函数直接为

\[
\phi(k)=\int_{-\infty}^{\infty} e^{ikx}\delta(x-\mu)\,dx=e^{ik\mu}.
\]

反过来，由傅里叶反演公式，

\[
\delta(x-\mu)=\frac{1}{2\pi}\int_{-\infty}^{\infty} e^{-ik\mu}e^{ikx}\,dk
=\frac1{2\pi}\int_{-\infty}^{\infty} e^{ik(x-\mu)}\,dk .
\]

这就是 \(\delta\) 函数的积分表示。
\end{solution}
\end{question}

